\chapter{Pipeline: Coleta, Engenharia de Features e Modelos Iniciais}\label{cap:pipeline}
% ============================================================================

As três primeiras etapas do \textit{pipeline} são detalhadas a seguir: coleta de notícias financeiras, engenharia de \textit{features} (mercado e texto) e treinamento inicial de modelos com \textit{embeddings} genéricos.

\section{Coleta de notícias financeiras}\label{sec:coleta}

O corpus textual foi construído a partir do portal InfoMoney, uma das principais fontes de notícias financeiras no Brasil. O módulo \texttt{extractor.py} acessa a REST API do WordPress do portal (\texttt{/wp-json/wp/v2/posts}), realizando buscas paginadas (100 artigos por página) para cada \textit{ticker} analisado. Cada requisição HTTP possui mecanismo de retentativa: até 3 tentativas com \textit{backoff} linear (1\,s, 2\,s), tratando erros de \textit{timeout}, conexão e códigos de status 429/5xx. Os artigos passam por pré-processamento (remoção de \textit{tags} HTML, decodificação de entidades, normalização de \textit{whitespace}) e a extração ocorre em paralelo via \texttt{ThreadPoolExecutor}.

\begin{table}[htb]
\centering
\caption{Volume de artigos coletados por ativo.}\label{tab:artigos}
\begin{tabular}{lrr}
\toprule
\textbf{Ativo} & \textbf{Artigos} & \textbf{Período} \\
\midrule
ITUB4 (Itaú Unibanco) & 2.572 & 2009--2026 \\
PETR4 (Petrobras)     & 1.775 & 2009--2026 \\
VALE3 (Vale)          & 1.525 & 2009--2026 \\
\midrule
\textbf{Total}        & \textbf{5.872} & \\
\bottomrule
\end{tabular}
\end{table}

\section{Engenharia de features}\label{sec:features}

O módulo \texttt{yahoo\_finance.py} obtém dados OHLCV via \texttt{yfinance} e calcula 11 \textit{features} técnicas por dia: fechamento, volume, retorno diário, médias móveis (7 e 21 dias), volatilidade (desvio-padrão 21 dias) e 5 valores defasados. O módulo \texttt{news\_embedder.py} transforma cada artigo em um vetor de 1.024 dimensões via Ollama (\texttt{qwen3-embedding:4b}), com agregação diária por média simples. As \textit{features} são combinadas via \textit{left join} por data com \textit{forward-fill} (\textit{forward-fill}: propagação do último valor conhecido para dias sem publicação de notícias). O \textit{dataset} resultante contém 1.227 dias $\times$ 1.035 \textit{features} (11 de preço + 1.024 de \textit{embeddings}).

\section{Modelos iniciais com embeddings genéricos}\label{sec:modelos_iniciais}

O alvo binário é definido como 1 se $\text{Close}[t+21] > \text{Close}[t]$ (horizonte $h=21$ dias), com desbalanceamento de 59\% ``Sobe'' / 41\% ``Desce''. Este horizonte é mantido nos Capítulos~\ref{cap:pipeline} e~\ref{cap:sentimento}; o Capítulo~\ref{cap:investigacao} também explora $h=5$ dias na varredura de horizontes (\texttt{horizon\_sweep.ipynb}). Os \textit{embeddings} são reduzidos para 32 componentes via PCA (\textit{Principal Component Analysis}), ajustado exclusivamente no conjunto de treino para evitar vazamento de dados. A escolha de 32 componentes foi baseada na inspeção da curva de variância explicada acumulada, conforme reportado em \texttt{3.model\_traning/feature\_engineering.ipynb}: os primeiros 32 componentes retêm 61,4\% da variância total dos \textit{embeddings} do conjunto de treino. A validação segue \textit{split} cronológico 70/15/15. Neste protocolo, os primeiros 70\% dos dias formam o conjunto de treino, os 15\% seguintes o conjunto de validação, e os 15\% finais o conjunto de teste, mantendo a ordem cronológica e evitando vazamento de dados futuros.

\begin{table}[htb]
\centering
\caption{Resultados com \textit{embeddings} genéricos Ollama (Etapa~3). \textbf{Nota:} estimativas pontuais de uma única semente e uma única janela temporal, reportadas para continuidade narrativa; o Capítulo~\ref{cap:investigacao} aplica avaliação multi-semente e multi-\textit{fold} aos resultados da Etapa~4.}\label{tab:etapa3}
\begin{tabular}{lcc}
\toprule
\textbf{Modelo} & \textbf{Arquitetura} & \textbf{ROC-AUC} \\
\midrule
BiLSTM Original  & 2 camadas, 128 \textit{hidden}, 30\% \textit{dropout} & 0,443 \\
BiLSTM Reduzido  & 1 camada, 32 \textit{hidden}, 50\% \textit{dropout}  & 0,505 \\
Transformer      & 2 camadas, 4 cabeças, $d_\text{model}=64$              & 0,568 \\
\textbf{XGBoost} & 300 árvores, profundidade 4                            & \textbf{0,610} \\
\bottomrule
\end{tabular}
\end{table}

Todos os modelos apresentaram desempenho próximo do acaso, sugerindo que \textit{embeddings} genéricos de alta dimensionalidade não são adequados para esta tarefa. A etapa seguinte explora uma representação alternativa, compacta e específica ao domínio (Capítulo~\ref{cap:sentimento}).


% ============================================================================
